% ============================ 4. TECHNOLOGY BETS ==================================
\newpage
\phase{Section four}{The less-chosen technologies}{Seven bets that most of the market is not making yet. For each: how ready it is, which way costs are heading, who is funding it, and what has to happen for it to win.}

\subsection{How to read readiness}

Technology readiness level (TRL) runs from 1 (an idea on paper) to 9 (proven in commercial operation). The TRLs below are \textbf{my estimates}, based on the most advanced public deployment I could find, so every one is \estd.

{\tablefont
\begin{xltabular}{\textwidth}{@{}>{\RaggedRight\arraybackslash}p{19mm} c L L L L@{}}
\caption{Seven technology bets.}\\
\toprule
\thead{Bet} & \thead{TRL} & \thead{Evidence and cost direction} & \thead{Who is backing it} & \thead{What must happen to win} & \thead{Fit for an Indian entrant}\\ \midrule
\endfirsthead
\toprule
\thead{Bet} & \thead{TRL} & \thead{Evidence and cost direction} & \thead{Who is backing it} & \thead{What must happen to win} & \thead{Fit for an Indian entrant}\\ \midrule
\endhead
\textbf{AEM electrolysis} & 6–7 & First 5 MW system delivered (April 2026) \cite{enki2026}. Enapter took €4.5 m of returned-stack costs in 2025 \cite{enapter2026} \rep. It uses nickel catalysts rather than precious metals \cite{irena2020}, so costs should fall with scale \estd & Enapter, Versogen (licensing a 250–300 MW Indian factory) \cite{enki2026} & 20,000+ hours of field durability at multi-MW scale; makers that survive long enough to honour warranties & \score{3} Good as a supply-chain play; risky as a plant foundation\\
\textbf{SOEC} & 6–7 & IEA capex range \usd{}2,800–5,600/kW \cite{h2techcapex2022}. Leader Topsoe ``hibernated'' its factory and cancelled a US plant (June 2026) \cite{h2central2026} \rep & Topsoe, Sunfire, Bloom Energy & Stack life far beyond 20,000 h; cheap industrial waste heat next door & \score{1} Wrong fit for solar-led India today\\
\textbf{Low-iridium PEM} & 8–9 & ITM cut iridium a further 40\% (Nov 2024) on top of earlier cuts over 80\% \cite{itm2024}. World iridium supply is only 7–8 t a year \cite{h2tech2022} \rep & ITM, Johnson Matthey, Electric Hydrogen & Thrifting plus 90\%+ recycling of iridium at end of life \cite{h2tech2022} & \score{3} Catalyst recycling is an Indian niche worth a look\\
\textbf{Direct seawater electrolysis} & 4–5 & A membrane-based seawater electrolyser ran for over 3,200 hours in lab conditions \cite{xie2022} \rep\ \confM. But desalinating water first costs at most \usd{}0.10/kg of hydrogen \cite{rmi2023} & Chinese academic and state groups & A cost saving larger than \usd{}0.10/kg, which is hard; mainly useful offshore & \score{1} Solves a problem that cheap desalination already solves\\
\textbf{Ammonia cracking} & 6–7 at scale & Amogy: nearly \usd{}300 m raised, a \usd{}700 m valuation, an ammonia tugboat in 2024, a 40 MW power plant targeted for 2028–29 \cite{amogy2025} \rep\prj & Amogy, Topsoe, KBR & Large import terminals in Europe and Japan & \score{2} India exports ammonia; cracking happens at the buyer's end\\
\textbf{Natural (white) hydrogen} & 3–5 & Koloma raised \usd{}245 m (Series B) to find underground hydrogen; deposits so far found mostly by accident, as in Mali \cite{koloma} \rep. The IEA warns deposits may be too scattered to pay \cite{koloma} & Koloma (Khosla, Breakthrough Energy Ventures) & Proof of large, repeatable, cheap-to-tap reservoirs & \score{1} Exploration capital and geology data an MBA start-up lacks\\
\textbf{Biomass and waste gasification} & 5–7 & Counts as ``green'' under India's standard \cite{drishti2023}. NTPC NETRA has proposed a 1 t/day plasma-gasification plant on municipal and farm waste (Dec 2025) \cite{ntpcnetra2025} \prj & NTPC, Indian research institutes & Cheap, reliable feedstock collection; meeting the 2 kg CO\textsubscript{2} limit with certification & \score{2} Interesting for small, distributed plants near waste sources\\
\bottomrule
\end{xltabular}}

\subsection{Patents and money: what the funding trail says}

The best study of hydrogen patents, by the European Patent Office and the IEA, covers 2011–2020. The EU held 28\% of international hydrogen patent families, Japan 24\% and the US 20\%. Electrolysis drove the growth in clean-production patents. The finding that matters most for a founder is about money: \textbf{start-ups that held patents received over half of the roughly \usd{}10 bn of venture capital invested in hydrogen firms, and over 80\% of the late-stage money}, even though they were under a third of the start-ups \cite{epo2023} \rep. \confH

\begin{finding}[Key finding \textbullet\ in a technology bet, the patent is the ticket]
If we ever make a technology bet, the investment case starts with defensible IP, not with a pilot plant. That makes the realistic Indian technology plays \emph{component} bets (electrodes, catalysts, catalyst recycling, power electronics tuned for variable solar) rather than whole-stack bets. Newtrace has already chosen this route \cite{newtrace2026}. \confM
\end{finding}

\begin{verdict}[My ranking of the seven bets for us]
\textbf{AEM} and \textbf{low-iridium PEM with catalyst recycling} are the only two worth tracking as a business, and both through components and services, not plants. Seawater electrolysis and ammonia cracking solve problems India does not have. SOEC and natural hydrogen need capital, heat or geology a start-up cannot reach. Biomass is a small, local niche: worth a pilot, not a strategy. \confM
\end{verdict}
